\documentclass[11pt]{article}
\usepackage[margin=1in]{geometry}
\usepackage{amsmath,amssymb,amsthm}
\usepackage{booktabs}
\usepackage{graphicx}
\usepackage{float}
\usepackage{enumitem}
\usepackage[hidelinks]{hyperref}

% ---------- theorem environments (numbered within section) ----------
\newtheorem{theorem}{Theorem}[section]
\newtheorem{proposition}[theorem]{Proposition}
\theoremstyle{definition}
\newtheorem{definition}[theorem]{Definition}
\newtheorem{assumption}{Assumption}

% ---------- macros ----------
\newcommand{\R}{\mathbb{R}}
\newcommand{\T}{^{\top}}
\newcommand{\tr}{\operatorname{tr}}
\newcommand{\E}{\mathbb{E}}
\newcommand{\diag}{\operatorname{diag}}
\newcommand{\rul}[1]{\text{[#1]}}
\newcommand{\owner}[2]{\par\medskip\noindent\textbf{Owner:} #1 \quad$|$\quad \textbf{Checked by:} #2\par\medskip}
% Includes a figure if the file exists; otherwise draws a labelled box so the document still compiles.
% Looks in figures/ (Overleaf) and ../../figures/ (repo: docs/report_iteration3/main.tex).
\newcommand{\fig}[2]{\IfFileExists{figures/#2}{\includegraphics[width=#1]{figures/#2}}{\IfFileExists{#2}{\includegraphics[width=#1]{#2}}{\IfFileExists{Figures/#2}{\includegraphics[width=#1]{Figures/#2}}{\IfFileExists{../../figures/#2}{\includegraphics[width=#1]{../../figures/#2}}{\fbox{\parbox[c][4cm][c]{#1}{\centering\small missing file: \texttt{\detokenize{#2}}}}}}}}}

\sloppy

\title{DS5020 Final Project: Iteration 3\\[4pt]
\large Anomaly Detection with the Multivariate Normal Distribution\\
Small Worked Example, From-Scratch Implementation and Verification}
\author{Bijun Wu \and Kanchan Patkar}
\date{November 2026}

\begin{document}
\maketitle

\begin{center}
GitHub: \url{https://github.com/WUB8/ds5020-anomaly-detection}\\[2pt]
Project Leader (Iteration 3): Kanchan Patkar
\end{center}

%=====================================================================
\section{Task table}
%=====================================================================
\begin{center}
\begin{tabular}{@{}lp{8.6cm}lll@{}}
\toprule
ID & Task & Owner & Checker & Status \\
\midrule
I1 & Small worked example, exact fractions (\texttt{small\_example.py}) & Kanchan & Bijun & Done \\
I2 & From-scratch methods (\texttt{methods.py}): MLE, Jacobi, LDL/Cholesky & Bijun & Kanchan & Done \\
I3 & From-scratch thresholds ($\chi^2$ and $F$ quantiles) & Kanchan & Bijun & Done \\
I4 & Simulation as functions (\texttt{simulate.py}) & Kanchan & Bijun & Done \\
I5 & Simulation test (\texttt{check\_simulation.py}) & Bijun & Kanchan & Done \\
I6 & Real data run and library comparison (\texttt{run\_real.py}) & Bijun & Kanchan & Done \\
I7 & Assumption checks (\texttt{diagnostics.py}) & Kanchan & Bijun & Done \\
I8 & Verification table & Kanchan & Bijun & Done \\
IQ1, IQ3, IQ4 & Math questions 1, 3, 4 & Kanchan & Bijun & Done \\
IQ2, IQ5 & Math questions 2, 5 & Bijun & Kanchan & Done \\
\bottomrule
\end{tabular}
\end{center}

%=====================================================================
\section{Background from Iteration 2}\label{sec:background}
%=====================================================================
This section recalls, without proofs, the model, results, algorithm and assumptions derived in the Iteration 2 report, so that this report can be read on its own. The proofs are in the Iteration 2 report (\texttt{docs/report\_iteration2.pdf}).

\subsection{Problem and data}
We model one genuine user's typing rhythm with a multivariate normal distribution and flag a new typing as an impostor when its squared Mahalanobis distance from the user's mean is too large. The data are the CMU Keystroke Dynamics Benchmark (Killourhy and Maxion, 2009) \cite{km09}: 51 subjects each typed the password \texttt{.tie5Roanl} + Return 400 times (8 sessions of 50 repetitions on different days). Each typing gives 11 hold times $H_k$, 10 keydown--keydown latencies $DD_{k,k+1}$ and 10 keyup--keydown latencies $UD_{k,k+1}=DD_{k,k+1}-H_k$. We use the $d=21$ features H and DD, in log seconds.

\subsection{Notation}
\begin{center}
\small
\begin{tabular}{@{}lll@{}}
\toprule
Symbol & Dimension & Meaning \\
\midrule
$n,\ d$ & scalars & training typings of the genuine user ($n=200$), features ($d=21$) \\
$X$ & $\R^{n\times d}$ & training data; row $i$ is $x_i\T$, $x_i=\log x_i^{\text{raw}}$ (entrywise) \\
$\mu,\ \Sigma$ & $\R^d,\ \R^{d\times d}$ & true mean and covariance (symmetric positive definite) \\
$\hat\mu,\ \hat\Sigma$ & $\R^d,\ \R^{d\times d}$ & maximum-likelihood estimates \\
$X_c$ & $\R^{n\times d}$ & centred data, $X_c=X-\mathbf 1\hat\mu\T$ \\
$Q,\ \Lambda$ & $\R^{d\times d}$ & eigenvectors and eigenvalues, $\hat\Sigma=Q\Lambda Q\T$, $\Lambda=\diag(\lambda_1,\dots,\lambda_d)$ \\
$W$ & $\R^{d\times d}$ & whitening matrix $W=\hat\Sigma^{-1/2}=Q\Lambda^{-1/2}Q\T$ \\
$v,\ z$ & $\R^d$ & deviation $v=x-\hat\mu$ and whitened vector \\
$d^2(x)$ & scalar & squared Mahalanobis distance $v\T\hat\Sigma^{-1}v$ (not $d$ squared) \\
$\alpha,\ \tau$ & scalars & target false-alarm rate and threshold $\tau=\chi^2_{d,1-\alpha}$ \\
$L,\ D$ & $\R^{d\times d}$ & $\hat\Sigma=LDL\T$: $L$ unit lower triangular, $D=\diag(D_1,\dots,D_d)$ (new) \\
$L_c$ & $\R^{d\times d}$ & Cholesky factor, $\hat\Sigma=L_cL_c\T$, $L_c=LD^{1/2}$ (new) \\
$w$ & $\R^d$ & solution of $Lw=v$ (new) \\
$\tau_F$ & scalar & small-sample threshold $\frac{d(n+1)}{n-d}F_{d,n-d,1-\alpha}$ (new) \\
$\operatorname{cond}(A)$ & scalar & condition number $\lambda_{\max}/\lambda_{\min}$ of a symmetric positive definite $A$ \\
$m_{\text{gen}},\ m_{\text{imp}}$ & scalars & genuine (200) and impostor (2500) test typings per user \\
$\varepsilon$ & scalar & ridge added to $\hat\Sigma$ if it is ill-conditioned \\
\bottomrule
\end{tabular}
\end{center}

\subsection{Results used (proved in Iteration 2)}
\begin{theorem}[MLE]\label{thm:mle}
For $x_1,\dots,x_n$ i.i.d.\ $\mathcal N(\mu,\Sigma)$, the log-likelihood is maximised by $\hat\mu=\frac1nX\T\mathbf 1$ and $\hat\Sigma=\frac1nX_c\T X_c$ (division by $n$).
\end{theorem}
\begin{proposition}[Inverse square root]\label{prop:invsqrt}
If $\Sigma=Q\Lambda Q\T$ with $Q\T Q=I$ and all $\lambda_j>0$, then $\Sigma^{-1/2}=Q\Lambda^{-1/2}Q\T$ is symmetric, $(\Sigma^{-1/2})^2=\Sigma^{-1}$ and $\Sigma^{-1/2}\Sigma\Sigma^{-1/2}=I$. A symmetric matrix is positive definite exactly when all its eigenvalues are positive.
\end{proposition}
\begin{theorem}[Distribution of $d^2$]\label{thm:chi2}
If $x\sim\mathcal N(\mu,\Sigma)$ and $d^2$ uses the true $\mu,\Sigma$, then $z=\Sigma^{-1/2}(x-\mu)\sim\mathcal N(0,I_d)$ and $d^2(x)=z\T z\sim\chi^2_d$.
\end{theorem}
\begin{proposition}[False-alarm rate]\label{prop:far}
With $\tau=\chi^2_{d,1-\alpha}$, $\Pr(d^2(x)>\tau)=\alpha$ for a normal reading (true parameters). The observed rate over $m$ readings has standard error $\sqrt{\alpha(1-\alpha)/m}$.
\end{proposition}
\begin{proposition}[Mean training distance]\label{prop:meand}
With the MLEs, $\frac1n\sum_{i=1}^n(x_i-\hat\mu)\T\hat\Sigma^{-1}(x_i-\hat\mu)=d$ exactly, for any data with $\hat\Sigma$ invertible.
\end{proposition}
\noindent\textbf{Small-sample result.} With plug-in estimates from $n$ training points and $T^2=\frac{n-1}{n}d^2$ for a new observation, $\frac{n(n-d)}{d(n+1)(n-1)}T^2\sim F_{d,n-d}$ \cite{jw}.

\subsection{Algorithm}
\textbf{Fitting:} (0) $X=\log X^{\text{raw}}$; (1) $\hat\mu=\frac1nX\T\mathbf 1$; (2) $X_c=X-\mathbf 1\hat\mu\T$; (3) $\hat\Sigma=\frac1nX_c\T X_c$; (4) $\hat\Sigma=Q\Lambda Q\T$, check $\lambda_{\min}>0$; (5) $W=Q\Lambda^{-1/2}Q\T$; (6) $\tau=\chi^2_{d,1-\alpha}$.
\textbf{Scoring a new typing:} (7) $x=\log x^{\text{raw}}$, $v=x-\hat\mu$; (8) $z=Wv$; (9) $d^2=z\T z$; (10) flag $x$ if $d^2>\tau$.
Fitting costs $O(nd^2+d^3)$, with about $9d^3$ operations for the eigendecomposition in step 4; scoring costs $O(d^2)$ per typing.

\subsection{Assumptions}
\begin{assumption}[iid]\label{as:iid} The training typings are independent and all drawn from the same $\mathcal N(\mu,\Sigma)$ (no drift over time).\end{assumption}
\begin{assumption}[Normality]\label{as:normal} For a fixed user, the log timings are multivariate normal.\end{assumption}
\begin{assumption}[Positive definite]\label{as:pd} $\Sigma$ and $\hat\Sigma$ are positive definite; $\hat\Sigma$ is, exactly when $X_c$ has full column rank ($v\T\hat\Sigma v=\frac1n\|X_cv\|^2$).\end{assumption}
\begin{assumption}[Clean training data]\label{as:clean} The training typings are all genuine and typical.\end{assumption}
\begin{assumption}[Large sample]\label{as:large} $n\gg d$, so that $\hat\mu\approx\mu$ and $\hat\Sigma\approx\Sigma$.\end{assumption}

%=====================================================================
\section{From-scratch implementation}\label{sec:impl}
%=====================================================================
\textbf{What changed since Iteration 2.} The dataset, the simulation model ($d=5$, $\mu=0$, $\Sigma=BB\T+5I$, $n=5000$, $\alpha=0.01$) and the seed (\texttt{np.random.default\_rng(0)}) are unchanged; \texttt{simulate\_normal(d=5, n=5000, seed=0)} returns exactly the Iteration 2 data. Two things changed: (i) the impostor set now has 50 typings per other subject ($m_{\text{imp}}=2500$), as the Iteration 3 handout asks; (ii) the Iteration 2 script is kept unchanged as \texttt{src/iteration2\_check.py}, and the simulation model moved into functions in \texttt{src/simulate.py}.

\textbf{Rules.} The method in \texttt{src/methods.py} uses NumPy arrays, matrix products and elementwise operations only. \texttt{numpy.linalg}, SciPy and scikit-learn appear only in the scripts that check our answers. Every function starts with \texttt{my\_} and has the equation it computes written above it.

\begin{center}
\small
\begin{tabular}{@{}clp{6.3cm}l@{}}
\toprule
Step & Function & Equation & Source \\
\midrule
0 & \texttt{my\_log\_transform} & $X_{ij}=\log X^{\text{raw}}_{ij}$ & Algorithm \\
1--3 & \texttt{my\_mle} & $\hat\mu=\frac1nX\T\mathbf 1$, $X_c=X-\mathbf 1\hat\mu\T$, $\hat\Sigma=\frac1nX_c\T X_c$ & Thm.~\ref{thm:mle} \\
4 & \texttt{my\_jacobi\_eig} & $\hat\Sigma=Q\Lambda Q\T$ (Jacobi rotations) & Prop.~\ref{prop:jacobi} \\
5 & \texttt{my\_inv\_sqrt} & $W=Q\Lambda^{-1/2}Q\T$ & Prop.~\ref{prop:invsqrt} \\
8--9 & \texttt{my\_whiten\_distance2} & $z=Wv$, $d^2=z\T z$ & Thm.~\ref{thm:chi2} \\
4$'$ & \texttt{my\_ldl}, \texttt{my\_cholesky} & $\hat\Sigma=LDL\T$, $L_c=LD^{1/2}$ & Prop.~\ref{prop:ldl} \\
8$'$--9$'$ & \texttt{my\_forward\_sub}, \texttt{my\_mahalanobis2} & $Lw=v$, $d^2=\sum_jw_j^2/D_j$ ($=\|L_c^{-1}v\|^2$) & Prop.~\ref{prop:ldl} \\
--- & \texttt{my\_inverse} & $\hat\Sigma^{-1}=L^{-\top}D^{-1}L^{-1}$ & Prop.~\ref{prop:ldl} \\
6 & \texttt{my\_threshold} & $\tau=\chi^2_{d,1-\alpha}$ by bisection on the $\chi^2_d$ CDF & Eq.~\eqref{eq:chi2cdf} \\
6$'$ & \texttt{my\_threshold\_F} & $\tau_F=\frac{d(n+1)}{n-d}F_{d,n-d,1-\alpha}$ & Eq.~\eqref{eq:tauF} \\
10 & \texttt{my\_decide}, \texttt{my\_rates} & flag if $d^2>\tau$; false-alarm and detection rates & Prop.~\ref{prop:far} \\
\bottomrule
\end{tabular}
\end{center}

The same functions run the tiny example (with \texttt{fractions.Fraction} where no square root is needed), the simulation and the real data. Two routes compute $d^2$: the Iteration 2 route through $W=\hat\Sigma^{-1/2}$, and a Cholesky route that needs no eigenvalues. The scripts check that the two routes agree with each other and with \texttt{numpy.linalg.inv}.

\subsection{Step 4: the Jacobi eigenvalue method}
\begin{proposition}[One Jacobi rotation]\label{prop:jacobi}
Let $A$ be symmetric with $a_{pq}\neq0$ ($p<q$). Put $\theta=\dfrac{a_{qq}-a_{pp}}{2a_{pq}}$, $t=\dfrac{\operatorname{sign}(\theta)}{|\theta|+\sqrt{\theta^2+1}}$, $c=\dfrac1{\sqrt{t^2+1}}$, $s=tc$, and let $J$ be the identity except $J_{pp}=J_{qq}=c$, $J_{pq}=s$, $J_{qp}=-s$. Then $J$ is orthogonal, $A'=J\T AJ$ has the same eigenvalues as $A$, $a'_{pq}=0$, and the sum of squares of the off-diagonal entries falls by $2a_{pq}^2$.
\end{proposition}
\begin{proof}
\textbf{Orthogonal.} The columns $p$ and $q$ of $J$ are $ce_p-se_q$ and $se_p+ce_q$; both have length $\sqrt{c^2+s^2}=1$ and their dot product is $cs-sc=0$ [$c^2+s^2=c^2(1+t^2)=1$].

\textbf{Same eigenvalues.} $A'=J\T AJ=J^{-1}AJ$ is similar to $A$ [$J\T=J^{-1}$].

\textbf{Zero entry.}
\begin{align*}
a'_{pq} &= (ce_p-se_q)\T A\,(se_p+ce_q) && \rul{$a'_{pq}=(Je_p)\T A(Je_q)$}\\
&= (c^2-s^2)a_{pq}+cs(a_{pp}-a_{qq}) && \rul{expand, $a_{qp}=a_{pq}$}\\
&= c^2a_{pq}\big(1-t^2-2\theta t\big). && \rul{$s=tc$, definition of $\theta$}
\end{align*}
$t$ is the root of smaller size of $t^2+2\theta t-1=0$, namely $t=-\theta+\operatorname{sign}(\theta)\sqrt{\theta^2+1}=\operatorname{sign}(\theta)/(|\theta|+\sqrt{\theta^2+1})$, so $a'_{pq}=0$.

\textbf{Off-diagonal decrease.} An orthogonal similarity keeps $\|A\|_F$. Only rows and columns $p,q$ change, and the $2\times2$ block on $\{p,q\}$ keeps its Frobenius norm, so $(a'_{pp})^2+(a'_{qq})^2=a_{pp}^2+a_{qq}^2+2a_{pq}^2$. The diagonal therefore gains $2a_{pq}^2$ and the off-diagonal part loses it.
\end{proof}
\noindent Applying rotations to every pair $(p,q)$ in turn (one \emph{sweep}) and repeating drives the off-diagonal part to zero \cite{gvl}. The product of the rotations is $Q$ and the final diagonal is $\Lambda$, so $\hat\Sigma=Q\Lambda Q\T$. We stop when $\big(\sum_{j\neq k}a_{jk}^2\big)^{1/2}\le10^{-14}\|\hat\Sigma\|_F$ and sort the eigenvalues in increasing order.

\subsection{Steps 4$'$--9$'$: the $LDL\T$ and Cholesky route}
\begin{proposition}[$LDL\T$ factorization and $d^2$]\label{prop:ldl}
If $\hat\Sigma$ is symmetric positive definite, it can be written $\hat\Sigma=LDL\T$ with $L$ unit lower triangular and $D=\diag(D_1,\dots,D_d)$, $D_j>0$, computed column by column as
\[
D_j=\hat\Sigma_{jj}-\sum_{k<j}L_{jk}^2D_k,\qquad L_{ij}=\frac{1}{D_j}\Big(\hat\Sigma_{ij}-\sum_{k<j}L_{ik}L_{jk}D_k\Big)\quad(i>j).
\]
For $v=x-\hat\mu$, let $w$ solve $Lw=v$. Then
\[
d^2(x)=v\T\hat\Sigma^{-1}v=\sum_{j=1}^d\frac{w_j^2}{D_j}=\|z\|^2,\qquad z=L_c^{-1}v,\quad L_c=LD^{1/2},
\]
and $\hat\Sigma^{-1}=L^{-\top}D^{-1}L^{-1}$.
\end{proposition}
\begin{proof}
\textbf{Formulas.} Entry $(i,j)$, $i\ge j$, of $LDL\T$ is $\sum_{k\le j}L_{ik}D_kL_{jk}$ [$L$ lower triangular]. With $L_{jj}=1$, setting it equal to $\hat\Sigma_{ij}$ and solving for the $k=j$ term gives the two formulas.

\textbf{Positive pivots.} Let $\hat\Sigma_{[j]}$ be the leading $j\times j$ block. Then $\hat\Sigma_{[j]}=L_{[j]}D_{[j]}L_{[j]}\T$ and $\det L_{[j]}=1$, so $\det\hat\Sigma_{[j]}=D_1\cdots D_j$ and $D_j=\det\hat\Sigma_{[j]}/\det\hat\Sigma_{[j-1]}$. Each $\hat\Sigma_{[j]}$ is positive definite (take $v$ with zeros after entry $j$ in $v\T\hat\Sigma v>0$), so its determinant, the product of its positive eigenvalues, is positive. Hence $D_j>0$, and the algorithm never divides by zero. Conversely, a pivot $D_j\le0$ shows that $\hat\Sigma$ is \emph{not} positive definite; \texttt{my\_ldl} stops with an error in that case (Assumption~\ref{as:pd}).

\textbf{Inverse.} $\hat\Sigma^{-1}=(LDL\T)^{-1}=L^{-\top}D^{-1}L^{-1}$ [$(ABC)^{-1}=C^{-1}B^{-1}A^{-1}$]; $L$ is invertible because $\det L=1$.

\textbf{Distance.}
\begin{align*}
v\T\hat\Sigma^{-1}v &= v\T L^{-\top}D^{-1}L^{-1}v=(L^{-1}v)\T D^{-1}(L^{-1}v) && \rul{$(L^{-1}v)\T=v\T L^{-\top}$}\\
&= w\T D^{-1}w=\sum_j w_j^2/D_j && \rul{$w=L^{-1}v$, $D$ diagonal}\\
&= \|D^{-1/2}w\|^2=\|L_c^{-1}v\|^2. && \rul{$L_c^{-1}=D^{-1/2}L^{-1}$}
\end{align*}
\end{proof}
\noindent\textbf{Why both routes give the same $d^2$.} $z_W=Wv$ and $z_L=L_c^{-1}v$ are both whitened vectors. They differ by the matrix $U=WL_c$, which is orthogonal: $UU\T=WL_cL_c\T W=W\hat\Sigma W=I$ (Prop.~\ref{prop:invsqrt}). An orthogonal matrix keeps lengths, so $\|z_W\|=\|z_L\|$.

\noindent\textbf{Cost.} The $LDL\T$ factorization takes about $d^3/3$ operations and each forward substitution about $d^2$, against about $9d^3$ for the eigendecomposition (Section~\ref{sec:background}, step 4). Forward substitution $w_i=v_i-\sum_{k<i}L_{ik}w_k$ uses only additions, multiplications and divisions, so with fractions it is exact.

\subsection{Step 6: the thresholds without SciPy}
The CDF of $\chi^2_d$ is the regularized lower incomplete gamma function with $a=d/2$, $x=t/2$:
\begin{equation}\label{eq:chi2cdf}
F_{\chi^2_d}(t)=P(a,x)=\frac{e^{-x}x^a}{\Gamma(a)}\sum_{k=0}^\infty\frac{x^k}{a(a+1)\cdots(a+k)},
\end{equation}
where $\Gamma(a)$ for $a\in\{\tfrac12,1,\tfrac32,\dots\}$ follows from $\Gamma(1)=1$, $\Gamma(\tfrac12)=\sqrt\pi$ and $\Gamma(z+1)=z\Gamma(z)$. All terms are positive, and we stop adding when a term is below $10^{-17}$ of the sum. $F_{\chi^2_d}$ is continuous and increasing, so $\tau=\chi^2_{d,1-\alpha}$ is found by bisection on $F_{\chi^2_d}(t)=1-\alpha$. For $d=2$, $F_{\chi^2_2}(t)=1-e^{-t/2}$ and $\tau=-2\ln\alpha$ exactly, which we use as a check.

For the small-sample threshold (Assumption~\ref{as:large}), the Iteration 2 result $\frac{n(n-d)}{d(n+1)(n-1)}T^2\sim F_{d,n-d}$ with $T^2=\frac{n-1}{n}d^2$ gives
\begin{equation}\label{eq:tauF}
\frac{n-d}{d(n+1)}\,d^2(x)\sim F_{d,n-d}\quad\Longrightarrow\quad \tau_F=\frac{d(n+1)}{n-d}\,F_{d,n-d,1-\alpha}
\end{equation}
for a new typing $x$. The $F$ CDF is a regularized incomplete beta function, $F_{F_{d_1,d_2}}(y)=I_{d_1y/(d_1y+d_2)}(\tfrac{d_1}2,\tfrac{d_2}2)$, which we evaluate with its continued fraction \cite{nr} and invert by bisection. For $n=200$, $d=21$, $\alpha=0.01$: $\tau=38.932$ and $\tau_F=46.183$.

%=====================================================================
\section{Small worked example}\label{sec:small}
%=====================================================================
Same data as the Iteration 2 hand-worked example: $x_1=(1,2)$, $x_2=(3,4)$, $x_3=(5,4)$, $x_4=(3,6)$, new points $(5,6)$ and $(5,2)$. \texttt{src/small\_example.py} runs the \texttt{my\_} functions on \texttt{Fraction} arrays and prints every intermediate result; all values below are copied from its output.

\textbf{Steps 1--3 (Thm.~\ref{thm:mle}).} Column sums $(12,16)$, $n=4$, so $\hat\mu=(3,4)$;
\[
X_c=\begin{pmatrix}-2&-2\\0&0\\2&0\\0&2\end{pmatrix},\qquad X_c\T X_c=\begin{pmatrix}8&4\\4&8\end{pmatrix},\qquad \hat\Sigma=\begin{pmatrix}2&1\\1&2\end{pmatrix}.
\]

\textbf{Step 4$'$ ($LDL\T$, Prop.~\ref{prop:ldl}).}
$D_1=\hat\Sigma_{11}=2$;\quad $L_{21}=\hat\Sigma_{21}/D_1=\tfrac12$;\quad $D_2=\hat\Sigma_{22}-L_{21}^2D_1=2-\tfrac12=\tfrac32$. So
\[
L=\begin{pmatrix}1&0\\\tfrac12&1\end{pmatrix},\qquad D=\diag\big(2,\tfrac32\big),\qquad LDL\T=\begin{pmatrix}2&1\\1&2\end{pmatrix}=\hat\Sigma\ \checkmark
\]
Both pivots are positive, so $\hat\Sigma$ is positive definite (Assumption~\ref{as:pd}).

\textbf{Inverse.} Forward substitution on the columns of $I$ gives $L^{-1}=\begin{pmatrix}1&0\\-\tfrac12&1\end{pmatrix}$ and
\[
\hat\Sigma^{-1}=L^{-\top}D^{-1}L^{-1}=\begin{pmatrix}\tfrac23&-\tfrac13\\-\tfrac13&\tfrac23\end{pmatrix}.
\]

\textbf{Steps 8$'$--9$'$ (distances).}
\begin{align*}
x=(5,6):&\quad v=(2,2),\quad w_1=2,\ w_2=2-\tfrac12\cdot2=1,\quad d^2=\tfrac{2^2}{2}+\tfrac{1^2}{3/2}=2+\tfrac23=\tfrac83,\\
x=(5,2):&\quad v=(2,-2),\quad w_1=2,\ w_2=-2-\tfrac12\cdot2=-3,\quad d^2=\tfrac{2^2}{2}+\tfrac{(-3)^2}{3/2}=2+6=8.
\end{align*}

\textbf{Steps 6 and 10.} With $\alpha=0.05$, \texttt{my\_threshold} gives $\tau=5.9914645471$, equal to $-2\ln0.05$ to $4\times10^{-15}$. Since $\tfrac83<\tau<8$, only $(5,2)$ is flagged.

\textbf{Second-way checks (exact).}
\begin{itemize}
\item $\hat\Sigma\hat\Sigma^{-1}=I$ exactly, and $\hat\Sigma^{-1}$ equals the $2\times2$ formula $\frac{1}{\det\hat\Sigma}\begin{pmatrix}2&-1\\-1&2\end{pmatrix}$ with $\det\hat\Sigma=3$.
\item $d^2$ computed directly as $v\T\hat\Sigma^{-1}v$ gives $\tfrac83$ and $8$, the same fractions as the $LDL\T$ route.
\item The training distances are $\tfrac83,0,\tfrac83,\tfrac83$ with mean exactly $2=d$ (Prop.~\ref{prop:meand}).
\end{itemize}

\textbf{Floating-point steps.} Square roots cannot be done in fractions, so steps 4, 5 and the Cholesky factor are computed in floating point and compared with the exact values: Jacobi gives eigenvalues $(1,3)$ (error $4\times10^{-16}$) and eigenvectors $\pm\tfrac1{\sqrt2}(1,-1)$, $\pm\tfrac1{\sqrt2}(1,1)$; $W^2$ matches the exact $\hat\Sigma^{-1}$ to $2\times10^{-16}$; $L_c=\begin{pmatrix}\sqrt2&0\\1/\sqrt2&\sqrt{3/2}\end{pmatrix}$ to $1\times10^{-16}$; and $d^2$ via $W$ and via $L_c$ match $(\tfrac83,8)$ to $4\times10^{-15}$ and $0$.

%=====================================================================
\section{Simulation test: recover the known truth}\label{sec:simtest}
%=====================================================================
Code: \texttt{src/check\_simulation.py}. Data: \texttt{simulate\_normal(d=5, n=5000, seed=0)}, the Iteration 2 data.

\subsection{Comparison with the library}
\begin{table}[H]
\centering
\small
\begin{tabular}{@{}lll@{}}
\toprule
Our function & Library (check only) & $\max|\text{diff}|$ \\
\midrule
\texttt{my\_mle}: $\hat\Sigma$ & \texttt{np.cov(X.T, bias=True)} & $1.8\times10^{-15}$ \\
\texttt{my\_cholesky} & \texttt{numpy.linalg.cholesky} & $4.4\times10^{-16}$ \\
\texttt{my\_inverse} & \texttt{numpy.linalg.inv} & $2.8\times10^{-17}$ \\
\texttt{my\_jacobi\_eig}: eigenvalues & \texttt{numpy.linalg.eigh} & $1.3\times10^{-14}$ \\
\texttt{my\_jacobi\_eig}: eigenvectors (signs aligned) & \texttt{numpy.linalg.eigh} & $9.1\times10^{-15}$ \\
$d^2$, Cholesky route & $v\T(\texttt{inv}\,\hat\Sigma)v$ & $1.1\times10^{-14}$ \\
$d^2$, $W$ route & $v\T(\texttt{inv}\,\hat\Sigma)v$ & $6.2\times10^{-14}$ \\
\texttt{my\_threshold}$(0.01,5)=15.0862724694$ & \texttt{scipy.stats.chi2.ppf} & $2.8\times10^{-14}$ \\
\bottomrule
\end{tabular}
\caption{From-scratch results against libraries on the simulated data. All differences are at rounding level: both sides compute the same formula on the same data, and the relative differences are a few times machine precision ($2.2\times10^{-16}$). Eigenvectors are only defined up to sign, so each of our eigenvectors is multiplied by the sign of its dot product with the library's before comparing.}
\label{tab:simlib}
\end{table}

\subsection{Comparison with the truth}
For normal data, $\operatorname{SE}(\hat\mu_j)=\sqrt{\Sigma_{jj}/n}$ and, to first order, $\operatorname{SE}(\hat\Sigma_{jk})=\sqrt{(\Sigma_{jk}^2+\Sigma_{jj}\Sigma_{kk})/n}$. Table~\ref{tab:simtruth} gives each error in standard errors, $z=(\text{estimate}-\text{truth})/\operatorname{SE}$.
\begin{table}[H]
\centering
\small
\begin{tabular}{@{}lrrrr@{}}
\toprule
Parameter & True & Estimate & SE & $z$ \\
\midrule
$\mu_1$ & 0.0000 & 0.0529 & 0.0339 & 1.56 \\
$\mu_2$ & 0.0000 & 0.0013 & 0.0443 & 0.03 \\
$\mu_3$ & 0.0000 & 0.0228 & 0.0498 & 0.46 \\
$\mu_4$ & 0.0000 & 0.0089 & 0.0379 & 0.24 \\
$\mu_5$ & 0.0000 & $-0.0150$ & 0.0407 & $-0.37$ \\
$\Sigma_{11}$ & 5.7413 & 5.7848 & 0.1148 & 0.38 \\
$\Sigma_{22}$ & 9.8247 & 9.8422 & 0.1965 & 0.09 \\
$\Sigma_{33}$ & 12.3961 & 12.2060 & 0.2479 & $-0.77$ \\
$\Sigma_{44}$ & 7.1888 & 7.1524 & 0.1438 & $-0.25$ \\
$\Sigma_{55}$ & 8.2660 & 8.1314 & 0.1653 & $-0.81$ \\
\bottomrule
\end{tabular}
\caption{Estimates against the truth ($n=5000$). The largest $|z|$ is 1.56 over $\mu$ and 1.99 over all 15 distinct entries of $\Sigma$, so every estimate is within 3 standard errors. $\|\hat\mu-\mu\|_2=0.060$ (typical size $\sqrt{\tr\Sigma/n}=0.093$) and $\|\hat\Sigma-\Sigma\|_F/\|\Sigma\|_F=0.029$.}
\label{tab:simtruth}
\end{table}

\subsection{Error against $n$}
For $n=50,100,200,500,1000,2000,5000,10000$ we keep the same truth and draw 200 independent data sets (\texttt{simulate\_normal(..., rep=r)}). Since $\hat\mu-\mu\sim\mathcal N(0,\Sigma/n)$, $\E\|\hat\mu-\mu\|^2=\tr\Sigma/n$, so the error should fall like $n^{-1/2}$.
\begin{figure}[H]
\centering
\begin{minipage}{0.52\textwidth}\centering\fig{\textwidth}{sim_error_vs_n.png}\end{minipage}\hfill
\begin{minipage}{0.45\textwidth}\centering\small
\begin{tabular}{@{}rccc@{}}
\toprule
$n$ & $\|\hat\mu-\mu\|$ & $\sqrt{\tr\Sigma/n}$ & rel.\ err.\ $\hat\Sigma$ \\
\midrule
50 & 0.934 & 0.932 & 0.322 \\
100 & 0.643 & 0.659 & 0.223 \\
200 & 0.457 & 0.466 & 0.159 \\
500 & 0.277 & 0.295 & 0.099 \\
1000 & 0.195 & 0.208 & 0.072 \\
2000 & 0.135 & 0.147 & 0.050 \\
5000 & 0.089 & 0.093 & 0.031 \\
10000 & 0.065 & 0.066 & 0.022 \\
\bottomrule
\end{tabular}
\end{minipage}
\caption{Mean error against the truth over 200 data sets per $n$, on log--log axes. The fitted slopes are $-0.507$ for $\hat\mu$ and $-0.503$ for $\hat\Sigma$, matching the theoretical $-1/2$. (The mean of $\|\hat\mu-\mu\|$ is slightly below $\sqrt{\tr\Sigma/n}=\sqrt{\E\|\hat\mu-\mu\|^2}$ by Jensen's inequality.)}
\label{fig:errn}
\end{figure}

\subsection{False-alarm rate and Q-Q plot}
We draw 5000 further points from the same model (rows 5001--10000 of \texttt{simulate\_normal(n=10000)}) and score them with the $\hat\mu,\hat\Sigma$ fitted on the first 5000. The standard error of a rate at $\alpha$ is $\sqrt{\alpha(1-\alpha)/5000}$.
\begin{table}[H]
\centering
\small
\begin{tabular}{@{}ccccc@{}}
\toprule
$\alpha$ & $\tau$ & Training & Held-out & $\alpha\pm3\,$SE \\
\midrule
0.10 & 9.2364 & 0.1010 & 0.1094 & $[0.0873,\,0.1127]$ \\
0.05 & 11.0705 & 0.0516 & 0.0554 & $[0.0408,\,0.0592]$ \\
0.01 & 15.0863 & 0.0102 & 0.0104 & $[0.0058,\,0.0142]$ \\
\bottomrule
\end{tabular}
\caption{False-alarm rates on simulated normal data; all lie within 3 standard errors of $\alpha$. At $\alpha=0.01$ the training rate 0.0102 is inside the pass band $0.010\pm0.004$.}
\label{tab:simfar}
\end{table}
A Kolmogorov--Smirnov test of the training $d^2$ against $\chi^2_5$ gives $p=0.47$, and the Q-Q plot of the held-out $d^2$ lies on the line (Figure~\ref{fig:stress}, left, blue).

\subsection{Stress tests}
\textbf{Heavy tails (Assumption~\ref{as:normal} broken).} \texttt{simulate\_t(d=5, n=10000, nu=3)} draws a multivariate $t_3$ with the same $\mu$ and scale matrix $\Sigma$; we fit on the first 5000 points and score the other 5000.
\begin{table}[H]
\centering
\small
\begin{tabular}{@{}ccc@{}}
\toprule
$\alpha$ & Normal data (held-out) & $t_3$ data (held-out) \\
\midrule
0.10 & 0.1094 & 0.1108 \\
0.05 & 0.0554 & 0.0890 \\
0.01 & 0.0104 & 0.0596 \\
\bottomrule
\end{tabular}
\caption{With heavy tails the false-alarm rate at $\alpha=0.01$ is 0.060, six times the target. The method is still computed correctly; the threshold is wrong because $d^2$ no longer follows $\chi^2_5$ (Thm.~\ref{thm:chi2} needs normality).}
\label{tab:t3}
\end{table}

\begin{figure}[H]
\centering
\fig{0.45\textwidth}{sim_iter3_d2_qq.png}\hfill
\fig{0.45\textwidth}{sim_stress_t3_qq.png}
\caption{Left: held-out $d^2$ on normal data against $\chi^2_5$ quantiles; the points lie on the line. Right: the same plot for $t_3$ data (red); the upper tail bends far above the line, which is why too many genuine points cross $\tau$.}
\label{fig:stress}
\end{figure}

\textbf{Duplicated feature (Assumption~\ref{as:pd} broken, IQ5).} Appending an exact copy of feature 1 makes column 6 of $X_c$ equal to column 1, so $X_c$ has rank 5 and $\hat\Sigma\in\R^{6\times6}$ is singular. Jacobi returns eigenvalues $(0.0000,\,5.24,\,7.52,\,9.63,\,11.65,\,14.87)$; \texttt{numpy.linalg.cond} gives $7.3\times10^{15}$, i.e.\ infinite up to rounding; and \texttt{my\_ldl} stops at pivot $D_6=2.7\times10^{-15}$, which is zero up to rounding (Prop.~\ref{prop:ldl}). A ridge $\hat\Sigma+\varepsilon I$ with $\varepsilon=10^{-3}\tr(\hat\Sigma)/d=0.0082$ brings the condition number down to $1.8\times10^3$; dropping the copy returns to the $d=5$ model with condition number 2.6. With an \emph{almost} collinear column (copy plus $\mathcal N(0,10^{-8})$ noise) the condition number is $3.0\times10^9$ and our $d^2$ differs from the \texttt{numpy.linalg.inv} result by a relative $2.3\times10^{-8}$: still correct, but about $10^8$ times further from the library than in the well-conditioned case, as expected when rounding errors are amplified by $\operatorname{cond}(\hat\Sigma)$.

\textbf{Small sample (Assumption~\ref{as:large}).} With $d=21$ and $n=200$ (the real-data sizes) we fit on 200 points and score 2000 new points, 200 times. Our thresholds match SciPy: $\tau=29.615,\,32.671,\,38.932$ and $\tau_F=34.238,\,38.090,\,46.183$ for $\alpha=0.10,0.05,0.01$.
\begin{table}[H]
\centering
\small
\begin{tabular}{@{}ccc@{}}
\toprule
$\alpha$ & $\chi^2_{21}$ threshold & $F$ threshold \eqref{eq:tauF} \\
\midrule
0.10 & 0.2124 & 0.1009 \\
0.05 & 0.1313 & 0.0504 \\
0.01 & 0.0432 & 0.0102 \\
\bottomrule
\end{tabular}
\caption{Mean held-out false-alarm rate over 200 repetitions with $n=200$, $d=21$. The plug-in $\chi^2$ threshold gives two to four times too many false alarms; the $F$ threshold gives $\alpha$. This reproduces Iteration 2, Part D, now with our own functions.}
\label{tab:smalln}
\end{table}

%=====================================================================
\section{Real data}\label{sec:real}
%=====================================================================
Code: \texttt{src/run\_real.py}.

\subsection{Cleaning steps}
\begin{enumerate}[leftmargin=*]
\item Load \texttt{DSL-StrongPasswordData.csv}: 20,400 rows $\times$ 34 columns, 51 subjects.
\item Missing values: 0; duplicate rows: 0. \textbf{No rows removed.}
\item Drop the 10 UD columns ($UD=DD-H$ exactly); keep $d=21$ features (11 H, 10 DD). All are $>0$, so the log is defined and no rows are removed.
\item Transform: we compare three versions: raw seconds, log of the 11 hold times only, and log of all 21 features (our model).
\item Split for each genuine user $u$: train = sessions 1--4 (200 typings), genuine test = sessions 5--8 (200), impostors = session 1 (reps 1--50) of each of the 50 other subjects (2500). No other train/test split is used.
\end{enumerate}

\subsection{Comparison with the library}
\begin{table}[H]
\centering
\small
\begin{tabular}{@{}lll@{}}
\toprule
Quantity (user \texttt{s002}, log features) & Library & max$|\text{diff}|$ (relative) \\
\midrule
$L_c$ (\texttt{my\_cholesky}) & \texttt{numpy.linalg.cholesky} & $1.1\times10^{-16}$ ($2\times10^{-16}$) \\
$\hat\Sigma^{-1}$ (\texttt{my\_inverse}) & \texttt{numpy.linalg.inv} & $3.6\times10^{-14}$ ($4\times10^{-16}$) \\
eigenvalues (\texttt{my\_jacobi\_eig}) & \texttt{numpy.linalg.eigh} & $9.4\times10^{-16}$ ($3\times10^{-15}$) \\
eigenvectors, signs aligned & \texttt{numpy.linalg.eigh} & $7.7\times10^{-15}$ ($7\times10^{-15}$) \\
$d^2$ of 2900 typings, Cholesky route & $v\T(\texttt{inv}\,\hat\Sigma)v$ & $4.6\times10^{-13}$ ($3\times10^{-16}$) \\
$d^2$ of 2900 typings, $W$ route & $v\T(\texttt{inv}\,\hat\Sigma)v$ & $3.2\times10^{-12}$ ($3\times10^{-15}$) \\
$\tau_F$, $\alpha=0.01$ & \texttt{scipy.stats.f.ppf} & $7\times10^{-13}$ \\
\midrule
$d^2$, all 51 users (genuine test) & $v\T(\texttt{inv}\,\hat\Sigma)v$ & relative $\le6.8\times10^{-16}$ \\
\bottomrule
\end{tabular}
\caption{Real data: our implementation against the libraries.}
\label{tab:reallib}
\end{table}
\noindent\textbf{Size of the differences.} The relative differences are $10^{-16}$--$10^{-15}$, a few units of machine precision, so they are rounding error. The absolute differences of $d^2$ ($\sim10^{-13}$) look larger only because the impostor distances are large (up to a few hundred). For \texttt{s002}, $\operatorname{cond}(\hat\Sigma)=67.9$, so rounding errors in solving with $\hat\Sigma$ are amplified by at most about $10^2$. The eigenvectors needed sign alignment: if $q$ is a unit eigenvector, so is $-q$, and the two programs need not choose the same sign. $d^2$ does not depend on the sign, because $W=\sum_j\lambda_j^{-1/2}q_jq_j\T$ is unchanged when $q_j\to-q_j$.

\subsection{Results}
\begin{table}[H]
\centering
\small
\begin{tabular}{@{}llcccc@{}}
\toprule
& & \multicolumn{2}{c}{$\alpha=0.01$} & \multicolumn{2}{c}{$\alpha=0.05$} \\
\cmidrule(lr){3-4}\cmidrule(l){5-6}
Features & Threshold & False alarm & Detection & False alarm & Detection \\
\midrule
raw seconds & $\chi^2_{21}$ & 0.340 & 0.875 & 0.390 & 0.926 \\
raw seconds & $F$ & 0.320 & 0.819 & 0.340 & 0.882 \\
log hold times only & $\chi^2_{21}$ & 0.350 & 0.900 & 0.405 & 0.941 \\
log hold times only & $F$ & 0.320 & 0.844 & 0.355 & 0.905 \\
log all 21 & $\chi^2_{21}$ & 0.445 & 0.960 & 0.530 & 0.980 \\
log all 21 & $F$ & 0.390 & 0.918 & 0.465 & 0.963 \\
\bottomrule
\end{tabular}
\caption{Genuine user \texttt{s002}: false-alarm rate on its own sessions 5--8 (200 typings) and detection rate on 2500 impostor typings.}
\label{tab:s002}
\end{table}

\begin{table}[H]
\centering
\small
\begin{tabular}{@{}llcccc@{}}
\toprule
& & \multicolumn{2}{c}{$\alpha=0.01$} & \multicolumn{2}{c}{$\alpha=0.05$} \\
\cmidrule(lr){3-4}\cmidrule(l){5-6}
Features & Threshold & False alarm & Detection & False alarm & Detection \\
\midrule
raw seconds & $\chi^2_{21}$ & 0.173 & 0.862 & 0.236 & 0.912 \\
raw seconds & $F$ & 0.127 & 0.799 & 0.178 & 0.868 \\
log hold times only & $\chi^2_{21}$ & 0.144 & 0.867 & 0.193 & 0.913 \\
log hold times only & $F$ & 0.111 & 0.810 & 0.149 & 0.874 \\
log all 21 & $\chi^2_{21}$ & 0.201 & 0.932 & 0.280 & 0.962 \\
log all 21 & $F$ & 0.141 & 0.889 & 0.209 & 0.936 \\
\bottomrule
\end{tabular}
\caption{Mean over all 51 genuine users. (Medians are similar: e.g.\ log all 21 with the $F$ threshold at $\alpha=0.01$ has median false-alarm rate 0.125 and median detection rate 0.924.)}
\label{tab:allusers}
\end{table}

\begin{figure}[H]
\centering
\fig{0.68\textwidth}{real_rates.png}
\caption{False-alarm and detection rates at $\alpha=0.01$, mean over the 51 users; the dashed line is the target $\alpha=0.01$.}
\label{fig:rates}
\end{figure}

\noindent\textbf{Reading the results.} (i) On the real data the false-alarm rate is 11--20\% instead of 1\%, for every transform and both thresholds. The simulation shows that the implementation is correct and that the $F$ threshold fixes the small-sample effect, so the gap must come from the data breaking the model's assumptions; Section~\ref{sec:assump} shows which ones. (ii) The log transform of all 21 features gives the highest detection rate (0.93 with $\chi^2$, 0.89 with $F$) but also the highest false-alarm rate. (iii) The $F$ threshold lowers both rates, because $\tau_F>\tau$; it trades about 4 percentage points of detection for about 6 points fewer false alarms. (iv) The detector separates genuine and impostor typings well (Figure~\ref{fig:realqq}, right), but the threshold from the normal theory is in the wrong place for this data.

%=====================================================================
\section{Verification table}\label{sec:verif}
%=====================================================================
Every number is printed by the script named in the first column (the scripts end with a block \texttt{VERIFICATION TABLE ROWS}).
\begin{table}[H]
\centering
\scriptsize
\begin{tabular}{@{}llllc@{}}
\toprule
Data & Quantity & Compared with & Difference & Pass? \\
\midrule
Small example (\texttt{small\_example.py}) & $\hat\Sigma\hat\Sigma^{-1}$ & $I$ (exact fractions) & 0 & \checkmark \\
Small example & $d^2$ ($LDL\T$) & $v\T\hat\Sigma^{-1}v$ (exact) & 0 & \checkmark \\
Small example & mean training $d^2$ & $d=2$ (exact) & 0 & \checkmark \\
Small example & $d^2$ via $W$ (float) & $(8/3,\,8)$ exact & $\max|\cdot|=4\times10^{-15}$ & \checkmark \\
Small example & $\tau$, $\alpha=0.05$ & $-2\ln\alpha$ & $4\times10^{-15}$ & \checkmark \\
\midrule
Simulation, $n=5000$ (\texttt{check\_simulation.py}) & $L_c$ & \texttt{numpy.linalg.cholesky} & $\max|\cdot|=4\times10^{-16}$ & \checkmark \\
Simulation, $n=5000$ & $\hat\Sigma^{-1}$ & \texttt{numpy.linalg.inv} & $\max|\cdot|=3\times10^{-17}$ & \checkmark \\
Simulation, $n=5000$ & eigenvalues & \texttt{numpy.linalg.eigh} & $\max|\cdot|=1\times10^{-14}$ & \checkmark \\
Simulation, $n=5000$ & $d^2$ (Cholesky) & \texttt{numpy.linalg.inv} & $\max|\cdot|=1\times10^{-14}$ & \checkmark \\
Simulation, $n=5000$ & $d^2$ ($W$) & \texttt{numpy.linalg.inv} & $\max|\cdot|=6\times10^{-14}$ & \checkmark \\
Simulation, $n=5000$ & $\tau=\chi^2_{5,0.99}$ & \texttt{scipy.stats.chi2} & $3\times10^{-14}$ & \checkmark \\
Simulation, $n=5000$ & $\hat\mu$ & true $\mu$ & $\|\cdot\|_2=0.060$, max $|z|=1.6$ & \checkmark \\
Simulation, $n=5000$ & $\hat\Sigma$ & true $\Sigma$ & rel.\ $\|\cdot\|_F=0.029$, max $|z|=2.0$ & \checkmark \\
Simulation, $n=50$--$10^4$ & error vs $n$ & slope $-1/2$ & $-0.51$ ($\hat\mu$), $-0.50$ ($\hat\Sigma$) & \checkmark \\
Simulation, $n=5000$ & false-alarm rate & $\alpha=0.01$ & $0.0102$ (band $0.010\pm0.004$) & \checkmark \\
Simulation (stress: $t_3$) & false-alarm rate & $\alpha=0.01$ & $0.060$: fails, as expected & \checkmark \\
Simulation (stress: copy) & pivot $D_6$ & 0 (singular) & cond $7\times10^{15}$, LDL stops & \checkmark \\
Simulation (stress: near copy) & $d^2$ & \texttt{numpy.linalg.inv} & rel.\ $2\times10^{-8}$ (cond $3\times10^{9}$) & \checkmark \\
Simulation (stress: $n=200$, $d=21$) & false-alarm rate & $\alpha=0.01$ & $\chi^2$: 0.043, $F$: 0.010 & \checkmark \\
\midrule
Real data, \texttt{s002} (\texttt{run\_real.py}) & $L_c$ & \texttt{numpy.linalg.cholesky} & rel.\ $2\times10^{-16}$ & \checkmark \\
Real data, \texttt{s002} & $\hat\Sigma^{-1}$ & \texttt{numpy.linalg.inv} & rel.\ $4\times10^{-16}$ & \checkmark \\
Real data, \texttt{s002} & eigenvalues & \texttt{numpy.linalg.eigh} & rel.\ $3\times10^{-15}$ & \checkmark \\
Real data, \texttt{s002} & $d^2$ (Cholesky) & \texttt{numpy.linalg.inv} & rel.\ $3\times10^{-16}$ & \checkmark \\
Real data, all 51 users & $d^2$ (Cholesky) & \texttt{numpy.linalg.inv} & max rel.\ $7\times10^{-16}$ & \checkmark \\
Real data, all users, log & false-alarm rate ($F$) & $\alpha=0.01$ & 0.141 (detection 0.889) & $\times$ \\
Real data, all users, raw & false-alarm rate ($F$) & $\alpha=0.01$ & 0.127 (detection 0.799) & $\times$ \\
\bottomrule
\end{tabular}
\caption{Verification table. For the stress tests, ``pass'' means the test shows the expected failure. The last two rows fail: the implementation agrees with the libraries on the real data, but the real data break Assumptions~\ref{as:iid} and~\ref{as:normal} (Section~\ref{sec:assump}).}
\label{tab:verif}
\end{table}

%=====================================================================
\section{Assumption checks}\label{sec:assump}
%=====================================================================
Each check runs first on simulated data, where the assumption holds by construction, to show what ``passing'' looks like, and then on the real data (user \texttt{s002}, log features; and over all 51 users). Code: \texttt{assumption\_checks} in \texttt{src/diagnostics.py}, called from both scripts.

\begin{table}[H]
\centering
\scriptsize
\begin{tabular}{@{}p{2.4cm}p{3.3cm}p{2.4cm}p{2.4cm}p{3.6cm}@{}}
\toprule
Assumption & Check & Simulation $d=5$, $n=5000$ & Simulation $d=21$, $n=200$ & Real data, \texttt{s002} (log) \\
\midrule
A1 iid (no drift) & Hotelling $T^2$, 1st vs 2nd half of training & $p=0.61$ & $p=0.018$ & $T^2=417$, $p=2\times10^{-33}$; rejects at 1\% for \textbf{51 of 51} users \\
& mean held-out $d^2$ (iid value $\frac{d(n+1)}{n-d-2}$) & 5.10 (5.01) & 24.1 (23.9) & 53.3 (23.9); median over users 31.0 \\
\midrule
A2 normality & features rejecting normality at 1\% (D'Agostino--Pearson) & 0 of 5 & 0 of 21 & 15 of 21 (log), 20 of 21 (raw); median over users 17 (log), 16 (raw) \\
& held-out $d^2$ vs exact law $d^2/c\sim F_{d,n-d}$ (KS) & $p=0.03$ & $p=0.54$ & $p=2\times10^{-33}$ \\
\midrule
A3 positive definite & $\lambda_{\min}$, $\operatorname{cond}(\hat\Sigma)$, smallest pivot & 5.05, 2.6, 5.78 & 14.3, 7.8, 26.2 & 0.0062, 67.9, 0.0119 \\
\midrule
A4 clean training & 5 largest training $d^2$ (vs $\chi^2_{d,0.999}$); change in $\hat\Sigma$ without them & 19.8--24.0 (20.5); 0.6\% & 33.3--38.3 (46.8); 8.3\% & 64.2--71.9 (46.8); 9.6\% \\
\midrule
A5 large sample & $n/d$ & 1000 & 9.5 & 9.5 \\
\bottomrule
\end{tabular}
\caption{Assumption checks. The simulation columns show what passing looks like; the real column shows which assumptions fail.}
\label{tab:assump}
\end{table}

\begin{figure}[H]
\centering
\fig{0.85\textwidth}{real_typing_time_by_rep.png}
\caption{Total typing time (sum of the 10 DD latencies, log scale) against repetition for four subjects, coloured by session, with session medians in red (from \texttt{src/load.py}, as in Iteration 2). Typing gets faster with practice, so sessions 5--8 do not follow the distribution fitted on sessions 1--4 (Assumption~\ref{as:iid}).}
\label{fig:practice}
\end{figure}

\noindent\textbf{Findings.}
\begin{itemize}
\item \textbf{A1 (iid) fails clearly.} Every one of the 51 users types differently in sessions 1--2 than in sessions 3--4 (practice effect, Figure~\ref{fig:practice}). For \texttt{s002} the genuine test typings have mean $d^2=53.3$, more than twice the 23.9 expected for new typings from the training distribution. This drift is the main reason the real false-alarm rate is above $\alpha$: the model is fitted to an earlier version of the user. (In the $n=200$ simulation, where there is no drift, $p=0.018$: one random draw that passes at the 1\% level.)
\item \textbf{A2 (normality) fails.} Most features reject normality even after the log transform. The log removes much of the skewness (median $|$skew$|$ for \texttt{s002} falls from 1.06 to 0.69), but heavy tails remain: the median excess kurtosis of the log features is 1.90 (max 12.7, for a normal distribution it is 0). Section~\ref{sec:iq} (IQ3, IQ4) gives the plots.
\item \textbf{A3 holds after dropping UD.} With the 21 log features $\hat\Sigma$ is well conditioned for every user (median condition number 52.8, max 214.7).
\item \textbf{A4 is not a large effect.} The 5 largest training distances of \texttt{s002} (64--72) exceed the $\chi^2_{21}$ 99.9\% point 46.8, a sign of heavy tails or unusual typings. Refitting without them changes $\hat\Sigma$ by 9.6\% and the $\alpha=0.01$ ($F$) rates from (0.390, 0.918) to (0.410, 0.934): almost no change.
\item \textbf{A5 is borderline but handled.} $n/d=9.5$; the $F$ threshold corrects the plug-in effect (Table~\ref{tab:smalln}), but this does not fix A1 or A2.
\end{itemize}

%=====================================================================
\section{Iteration 3 -- Math Questions}\label{sec:iq}
%=====================================================================
\subsection*{IQ1. Four points in $\R^2$: $\hat\mu$, $\hat\Sigma$, $\hat\Sigma^{-1}$ and the Mahalanobis distance}
Done step by step in Section~\ref{sec:small}, in exact fractions: $\hat\mu=(3,4)$, $\hat\Sigma=\begin{pmatrix}2&1\\1&2\end{pmatrix}$, $\hat\Sigma^{-1}=\tfrac13\begin{pmatrix}2&-1\\-1&2\end{pmatrix}$, and for the new point $(5,2)$, $v=(2,-2)$ and
\[
d^2=v\T\hat\Sigma^{-1}v=\tfrac13\big(2\cdot4-2\cdot(2)(-2)+2\cdot4\big)=\tfrac13(8+8+8)=8,
\]
while $(5,6)$ has $d^2=\tfrac83$. The second-way checks ($\hat\Sigma\hat\Sigma^{-1}=I$, the $LDL\T$ route, mean training $d^2=d$) all hold exactly.
\owner{Kanchan Patkar}{Bijun Wu}

\subsection*{IQ2. Inverse and Cholesky solve from scratch, compared with \texttt{numpy.linalg.inv}}
\texttt{my\_ldl} computes $\hat\Sigma=LDL\T$ (Prop.~\ref{prop:ldl}), \texttt{my\_cholesky} returns $L_c=LD^{1/2}$, \texttt{my\_forward\_sub} solves $Lw=v$, and \texttt{my\_inverse} forms $\hat\Sigma^{-1}=L^{-\top}D^{-1}L^{-1}$. Against NumPy the largest differences are $3\times10^{-17}$ (simulation) and $3.6\times10^{-14}$, relative $4\times10^{-16}$ (\texttt{s002}) for the inverse, and relative $\le7\times10^{-16}$ for $d^2$ over all 51 users (Tables~\ref{tab:simlib}, \ref{tab:reallib}). In practice we never form $\hat\Sigma^{-1}$ to score a typing: solving $Lw=v$ costs $d^2$ operations per typing and avoids the extra rounding of computing the inverse first.
\owner{Bijun Wu}{Kanchan Patkar}

\subsection*{IQ3. Q-Q plot of $d^2$ against $\chi^2_d$ on the real data}
\begin{figure}[H]
\centering
\fig{0.98\textwidth}{real_d2_qq.png}
\caption{User \texttt{s002}, log features. Left: sorted $d^2$ against $\chi^2_{21}$ quantiles for the training typings (blue) and the genuine test typings (green). Right: $d^2$ of genuine test typings and of 2500 impostor typings, with the $\chi^2_{21}$ density and the thresholds.}
\label{fig:realqq}
\end{figure}
\noindent\textbf{The fit is not good.} On the training typings the points follow the line in the middle but bend above it in the upper tail (largest $d^2=72$, where $\chi^2_{21}$ gives 44), a sign of heavy tails (compare the $t_3$ stress test, Figure~\ref{fig:stress}). On the genuine test typings the whole curve lies above the line, and from about the median on the test distances grow much faster than the $\chi^2_{21}$ quantiles, reaching three to six times them in the upper tail. This is the drift of Assumption~\ref{as:iid}, not only non-normality: the test typings come from later sessions. A KS test rejects both $\chi^2_{21}$ ($p=2\times10^{-44}$) and the exact small-sample law ($p=2\times10^{-33}$). The simulated Q-Q plot in Figure~\ref{fig:stress} shows what a good fit looks like.
\owner{Kanchan Patkar}{Bijun Wu}

\subsection*{IQ4. Normality of each feature}
\begin{figure}[H]
\centering
\fig{0.98\textwidth}{real_feature_qq_log.png}
\caption{Normal Q-Q plot of each log feature, user \texttt{s002}, 200 training typings (standardized). The red line is the normal reference. The raw-seconds version is \texttt{figures/real\_feature\_qq\_raw.png}; a simulated version, where all panels lie on the line, is \texttt{figures/sim\_feature\_qq.png}.}
\label{fig:featqq}
\end{figure}
\noindent Before the log, 20 of 21 features of \texttt{s002} reject normality at the 1\% level; the worst are the latencies \texttt{DD.i.e} (skewness 8.64), \texttt{DD.l.Return} (6.07), \texttt{DD.n.l} (4.31) and \texttt{DD.period.t} (3.02). After the log, 15 of 21 still reject. Most hold times are now close to normal (\texttt{H.period}, \texttt{H.i}, \texttt{H.e}, \texttt{H.five}, \texttt{H.a}, \texttt{H.l} lie near the line), while the DD latencies break the assumption most: \texttt{DD.l.Return} (skewness 2.87), \texttt{DD.n.l} (2.36), \texttt{DD.five.Shift.r} and \texttt{DD.Shift.r.o} have long right tails from hesitations, and \texttt{H.o} ($-2.20$) and \texttt{DD.a.n} ($-1.34$) have a few very short values on the left. Over all users the median number of rejecting features is 16 of 21 raw and 17 of 21 after the log: the log fixes skewness but not the heavy tails, and with $n=200$ the test detects even moderate excess kurtosis.
\owner{Kanchan Patkar}{Bijun Wu}

\subsection*{IQ5. What to do if $\hat\Sigma$ is singular or close to singular}
\textbf{Why it matters.} By Prop.~\ref{prop:ldl}, $\hat\Sigma$ is singular exactly when some pivot $D_j=0$, i.e.\ when the columns of $X_c$ are linearly dependent [$v\T\hat\Sigma v=\frac1n\|X_cv\|^2$]. Then $\hat\Sigma^{-1}$ does not exist and $d^2$ is undefined. If $\hat\Sigma$ is close to singular, $\lambda_{\min}$ is tiny and $d^2=\sum_j(q_j\T v)^2/\lambda_j$ divides by it, so noise in the direction $q_{\min}$ is blown up, and rounding errors grow like $\operatorname{cond}(\hat\Sigma)$ (stress test: relative error $2\times10^{-8}$ at $\operatorname{cond}=3\times10^9$ against $10^{-16}$ at $\operatorname{cond}=2.6$).

\textbf{Our data.}
\begin{center}
\small
\begin{tabular}{@{}lccc@{}}
\toprule
Features & $\operatorname{cond}(\hat\Sigma)$, \texttt{s002} & median over users & max over users \\
\midrule
31 raw (with UD) & $\infty$ (10 zero eigenvalues) & $\infty$ & $\infty$ \\
21 raw & $1.6\times10^3$ & 958 & 54{,}450 \\
21, log of hold times only & --- & 90 & 2{,}253 \\
21 log & 67.9 & 52.8 & 214.7 \\
\bottomrule
\end{tabular}
\end{center}
\textbf{What we do, in order.}
\begin{enumerate}[leftmargin=*]
\item \textbf{Remove exact dependences.} $UD=DD-H$, so we drop the 10 UD columns (Iteration 2). This turns a singular $\hat\Sigma$ into a positive definite one.
\item \textbf{Transform.} The log reduces the spread of scales between features (hesitations of several seconds versus holds of 0.1\,s); the median condition number falls from 958 to 53.
\item \textbf{Detect.} \texttt{my\_ldl} stops when a pivot $D_j\le10^{-12}\hat\Sigma_{jj}$, and we report $\operatorname{cond}(\hat\Sigma)$ for every user.
\item \textbf{Regularize if needed.} Use $\hat\Sigma+\varepsilon I$: every eigenvalue increases by $\varepsilon$, so $\operatorname{cond}=(\lambda_{\max}+\varepsilon)/(\lambda_{\min}+\varepsilon)$ is bounded. In the duplicated-feature stress test, $\varepsilon=0.0082$ turns $\operatorname{cond}=7\times10^{15}$ into $1.8\times10^3$. On our log features it is not needed: for \texttt{s002}, $\varepsilon=6\times10^{-5}$ changes the condition number from 67.9 to 67.3 and the $\alpha=0.01$ rates from (0.445, 0.960) to (0.445, 0.959). Other options are to keep only the eigen-directions with $\lambda_j$ above a cut-off (a pseudo-inverse), or to collect more training typings.
\end{enumerate}
\owner{Bijun Wu}{Kanchan Patkar}

%=====================================================================
\section{Member contributions}
%=====================================================================
\textbf{Kanchan Patkar} (Project Leader): worked the small example step by step and wrote \texttt{src/small\_example.py}; wrote the $\chi^2$ and $F$ threshold functions in \texttt{src/methods.py}, \texttt{src/simulate.py} and the assumption checks in \texttt{src/diagnostics.py}; wrote Sections~\ref{sec:small} and~\ref{sec:assump}, the verification table and the answers to IQ1, IQ3 and IQ4; checked Bijun's derivations (Props.~\ref{prop:jacobi}, \ref{prop:ldl}) and code. As leader, checked that each member did part of the math and part of the code.

\textbf{Bijun Wu}: derived the Jacobi and $LDL\T$ results (Props.~\ref{prop:jacobi}, \ref{prop:ldl}) and implemented them in \texttt{src/methods.py}; wrote \texttt{src/check\_simulation.py} and \texttt{src/run\_real.py}; wrote Sections~\ref{sec:impl}, \ref{sec:simtest} and~\ref{sec:real} and the answers to IQ2 and IQ5; checked Kanchan's worked example, threshold code and assumption checks.

%=====================================================================
\begin{thebibliography}{9}
\bibitem{km09} K.~S. Killourhy and R.~A. Maxion. Comparing anomaly-detection algorithms for keystroke dynamics. \emph{IEEE/IFIP International Conference on Dependable Systems and Networks (DSN)}, 2009.
\bibitem{jw} R.~A. Johnson and D.~W. Wichern. \emph{Applied Multivariate Statistical Analysis}, 6th ed. Pearson, 2007.
\bibitem{nr} W.~H. Press, S.~A. Teukolsky, W.~T. Vetterling and B.~P. Flannery. \emph{Numerical Recipes}, 3rd ed. Cambridge University Press, 2007.
\bibitem{gvl} G.~H. Golub and C.~F. Van Loan. \emph{Matrix Computations}, 4th ed. Johns Hopkins University Press, 2013.
\end{thebibliography}

\end{document}
